\documentclass[12pt]{article}
\usepackage[utf8]{inputenc}
\usepackage[spanish]{babel}
\usepackage{amsmath}
\usepackage{amsfonts}
\usepackage{amssymb}
\usepackage{geometry}
\geometry{letterpaper, margin=2.5cm}

\begin{document}

\begin{center}
    \Large \textbf{Evaluación de Examen 1: Unidad I} \\
    \large Física para Ciencias de la Salud (005-1713) \\
    \vspace{0.5cm}
    \textbf{Estudiante:} Daniela Escribano \quad \textbf{C.I.:} 32.031.568
    \vspace{0.3cm} \\
    \large \textbf{Calificación Final: 10/10 puntos}
\end{center}

\vspace{0.5cm}

\section*{Análisis de Resultados}

\subsection*{1. Ángulo plano (Conversión a radianes)}
\textbf{Procedimiento y resultado:} Plantea correctamente una regla de tres. Simplifica la fracción resultante $\left(\frac{75}{180}\right)$ dividiendo numerador y denominador entre 15 para obtener magistralmente el resultado exacto de $\frac{5\pi}{12}$ rad. Un procedimiento muy elegante. \\
\textbf{Veredicto:} \textbf{Correcto} (2/2 pts).

\subsection*{2. Sistemas de coordenadas (Longitud del hueso)}
\textbf{Procedimiento y resultado:} Utiliza la fórmula de distancia entre dos puntos y resuelve paso a paso las operaciones $\left(\sqrt{6^2 + 8^2} = \sqrt{100}\right)$. Aunque existe un pequeño error de transcripción en el recuadro final (escribe $\sqrt{10\text{cm}}$ en lugar de $10\text{cm}$), el procedimiento matemático es sólido y la gráfica adjunta corrobora la correcta comprensión geométrica. \\
\textbf{Veredicto:} \textbf{Correcto} (2/2 pts).

\subsection*{3. Vectores (Fuerza resultante)}
\textbf{Procedimiento y resultado:} Ejecuta una suma de vectores disponiendo las ecuaciones verticalmente. Esto le permite sumar las componentes $\hat{i}$ y $\hat{j}$ de forma clara, respetando las leyes de los signos y obteniendo el vector resultante correcto: $30\hat{i} + 60\hat{j}$ N. \\
\textbf{Veredicto:} \textbf{Correcto} (2/2 pts).

\subsection*{4. Potencias de diez (Notación científica y prefijo)}
\textbf{Procedimiento y resultado:} ¡Excelente resolución! La estudiante logra distinguir y cumplir con ambos requerimientos del problema: expresa la cifra en notación científica estricta ($1,25 \times 10^{-5}$ m) y, paralelamente, realiza el ajuste pertinente a $10^{-6}$ para nombrar de forma exacta el prefijo del S.I. ($12,5$ micrómetros). \\
\textbf{Veredicto:} \textbf{Correcto} (2/2 pts).

\subsection*{5. Sistemas de unidades (Conversión de densidad)}
\textbf{Procedimiento y resultado:} Presenta cierta confusión al intentar escribir la fórmula formal de los factores de conversión en cadena. Sin embargo, compensa esto desglosando las operaciones aritméticas paso a paso en la parte inferior: divide entre 1000 (para la masa) y multiplica por 1.000.000 (para el volumen), llegando a la respuesta final exacta de $1030 \text{ Kg/m}^3$. \\
\textbf{Veredicto:} \textbf{Correcto} (2/2 pts).

\vspace{0.5cm}
\section*{Comentarios Generales y Sugerencias}
¡Felicidades por obtener la calificación máxima! El examen evidencia un excelente razonamiento matemático, habilidades algebraicas muy sólidas y una interpretación lectora excepcional (destacando especialmente en la resolución de la Pregunta 4). 
\begin{itemize}
    \item Se sugiere practicar un poco más la notación formal de los factores de conversión en cadena (Pregunta 5) para que la presentación escrita acompañe la excelente lógica de cálculo demostrada.
\end{itemize}

\end{document}